\section*{الفصل \ref{ch:b2:magmatic-differentiation} --- التمايز الصهاري}

\begin{solution}{exo:b2:magmatic-differentiation:1}
(أ) خط سيولة An$_{50}$ هو \qty{1450}{\celsius}. وأول بلّورة تقع على
خط الصلابة عند تلك الدرجة، بين An$_{66.7}$ (عند 1372) وAn$_{83.3}$ (عند 1465):
$66.7 + 16.6\times(1450 - 1372)/(1465 - 1372) = 66.7 + 13.9 = \qty{80.6}{\percent}$
من الأنورثيت. (ب) عند \qty{1394}{\celsius} يكون السائل على خط السيولة،
An$_{33.3}$؛ والبلّورات على خط الصلابة، $66.7 + 16.6\times22/93 =
\qty{70.6}{\percent}$. وبقاعدة الرافعة يكون الكسر المتبلور $(50 -
33.3)/(70.6 - 33.3) = 16.7/37.3 = 0.448$. (ج) يختفي آخر سائل حين تبلغ
البلّورات An$_{50}$، عند درجة حرارة خط الصلابة \qty{1287}{\celsius}؛
وخط السيولة هناك يقع بين An$_{11.1}$ (عند 1265) وAn$_{20}$ (عند 1334): $11.1 +
8.9\times22/69 = \qty{13.9}{\percent}$ من الأنورثيت.
\end{solution}

\begin{solution}{exo:b2:magmatic-differentiation:2}
$\bar D_{\ce{Sr}} = 0.15\times0.000063 + 0.35\times0.157 + 0.50\times2.7 = 0.00001 +
0.0550 + 1.35 = 1.40$؛ $\bar D_{\ce{Zr}} = 0.0001 + 0.0683 + 0.0005 = 0.069$؛
$\bar D_{\ce{Rb}} = 0.00001 + 0.0012 + 0.0125 = 0.014$. السترونشيوم وحده، الذي يحمله
البلاجيوكليز، متوافق ($\bar D > 1$)؛ أمّا الزركونيوم والروبيديوم فغير متوافقين،
والروبيديوم أشدّ.
\end{solution}

\begin{solution}{exo:b2:magmatic-differentiation:3}
(أ) $C_L/C_0 = 0.5^{D-1}$: للسترونشيوم $0.5^{0.40} = 0.758$ (ينخفض)؛ وللزركونيوم $0.5^{-0.931} =
1.907$؛ وللروبيديوم $0.5^{-0.986} = 1.981$، قريبًا من الحدّ 2. (ب) تتغيّر Zr/Rb بالعامل
$0.5^{0.069 - 0.014} = 0.5^{0.055} = 0.963$، أي أقل من \qty{4}{\percent} بينما كاد كل
عنصر يتضاعف: فكما تقول النتيجة، تحتفظ نسبة عنصرين غير متوافقين
بقيمتها في \omterm{def:b2:magma-emplacement:magma}{الصهارة} الأمّ، فهي تخبر عن المصدر، لا
عن التبلور.
\end{solution}

\begin{solution}{exo:b2:magmatic-differentiation:4}
(أ) \ce{SiO2} $+5.03$، \ce{MgO} $-1.95$، \ce{K2O} $+0.82$ (في المئة)، والبوتاس
تضاعف بالعامل $1.67/0.85 = 1.96$. (ب) لعنصر غير متوافق تمامًا
$C_L/C_0 = 1/F$، ومنه $F = 0.85/1.67 = 0.509$: فكان لا بدّ أن يتبلور نصف
\omterm{def:b2:magma-emplacement:magma}{الصهارة}. (ج) إزالة \qty{12.5}{g} لكل \qty{100}{g} تترك $F \approx 0.88$
ولا ترفع البوتاس إلا بنحو $1/0.88 = 1.14$: فالتبلور وحده لا يستطيع
تفسير الارتفاع. لقد جاء البوتاس من الخارج، من الصخور الدخيلة الغرانيتية الغنية
بالبوتاس (\qty{3.6}{\percent})، التي استوعبتها \omterm{def:b2:magma-emplacement:magma}{الصهارة}.
\end{solution}

\begin{solution}{exo:b2:magmatic-differentiation:5}
(أ) $F = 0.8$، $D = 5$. التجزيئي: $0.8^{4} = 0.410$؛ المتوازن:
$1/(5 + 0.8\times(1 - 5)) = 1/1.8 = 0.556$. (ب) $F = 0.1$، $D = 0.01$. التجزيئي:
$0.1^{-0.99} = 9.77$؛ المتوازن: $1/(0.01 + 0.1\times0.99) = 1/0.109 = 9.17$.
(ج) حين $F \to 0$، يؤول قانون رايلي إلى $0$ من أجل $D > 1$ وإلى اللانهاية من أجل
$D < 1$، بينما يؤول قانون التوازن إلى $1/D$ (آخر سائل في توازن
مع بلّورات لها التركيب الكلي). \omterm{def:b2:magmatic-differentiation:crystallisation}{والتبلور التجزيئي} وحده يستطيع تجريد
سائل من عنصر متوافق تجريدًا تامًا.
\end{solution}

\begin{solution}{exo:b2:magmatic-differentiation:6}
(أ) تحوي وحدة الصيغة الواحدة 1.6 من \ce{MgO} و0.4 من \ce{FeO} وواحدًا من \ce{SiO2}: $1.6\times
40.30 = 64.48$، $0.4\times71.84 = 28.74$، 60.08، والمجموع \qty{153.30}{g}. فالنسب
المئوية الكتلية: \ce{MgO} \qty{42.06}{\percent}، \ce{FeO} \qty{18.75}{\percent}،
\ce{SiO2} \qty{39.19}{\percent}. (ب) قاعدة الرافعة على \ce{MgO}: $x = (10.0 - 8.0)/(42.06
- 8.0) = 2.0/34.06 = 0.0587$، أي نحو \qty{6}{\percent} من الأوليفين. (ج) سيليكا
السائل: $(49.0 - 0.0587\times39.19)/(1 - 0.0587) = 46.70/0.9413 =
\qty{49.61}{\percent}$. فإزالة الأوليفين، الأفقر بالسيليكا من البازلت، ترفع
سيليكا السائل، لكن ببطء فقط.
\end{solution}

\begin{solution}{exo:b2:magmatic-differentiation:7}
(أ) $F(0) = 0.768$؛ $F(1000) = 0.1091 - 0.5906 + 0.7678 = 0.286$؛ $F(2000) = 0.4364
- 1.1813 + 0.7678 = 0.023$. (ب) المميّز هو $(5.9064\times10^{-4})^2 -
4\times1.091\times10^{-7}\times0.7678 = 3.4885\times10^{-7} - 3.3507\times10^{-7} =
1.379\times10^{-8}$، وجذره التربيعي $1.174\times10^{-4}$؛ والجذر الأصغر هو
$H = (5.9064 - 1.174)\times10^{-4}/(2.182\times10^{-7}) = \qty{2169}{m}$، أي
أفق الساندويتش. (ج) $1 - 0.768 = 0.232$: فقد تبلور نحو ربع \omterm{def:b2:magma-emplacement:magma}{الصهارة}
قبل أن تتكوّن قاعدة المنطقة الدنيا المتكشّفة، في المنطقة المخفية
تحتها وفي السلسلة الحدّية على امتداد الجدران والسقف.
\end{solution}

\begin{solution}{exo:b2:magmatic-differentiation:8}
يحوي \omterm{def:b2:magmatic-differentiation:cumulate}{المتراكم} $D\,C_L = D\,C_0F^{D-1}$، فبين مستويين
$\ln(\mathrm{Ni}_2/\mathrm{Ni}_1) = (D - 1)\ln(F_2/F_1)$. (أ) $D - 1 = \ln(54/177)/
\ln(0.29/0.70) = (-1.187)/(-0.881) = 1.35$، ومنه $D \approx 2.35$ (والانحدار على
المتراكمات البالغ عددها 66 ذات $F \ge 0.23$ يعطي 2.38). (ب) $D - 1 = \ln(6.2/54)/\ln(0.17/0.29) = (-2.164)/(-0.534)
= 4.05$، ومنه $D \approx 5.05$. (ج) يتغيّر $D$ الكلي حين تتغيّر المجموعة
المتبلورة (معدن جديد يأخذ النيكل، أو معدن يغادر)،
ومعاملات التوزيع بين المعدن والمصهور نفسها ترتفع كلما برد السائل وتغيّر
تركيبه، كما قيس للأوليفين؛ والسائل المحبوس في المتراكمات يشوّش
التقدير.
\end{solution}

\begin{solution}{exo:b2:magmatic-differentiation:9}
(أ) $100 + 25.4 - 2.9 - 9.6 = \qty{112.9}{g}$. (ب) غرامات \ce{SiO2}: $55.53 +
25.4\times0.7248 - 2.9\times0.3919 - 9.6\times0.5054 = 55.53 + 18.41 - 1.14 - 4.85 =
\qty{67.95}{g}$، ومنه $67.95/112.9 = \qty{60.19}{\percent}$ مقابل 60.16. \ce{MgO}:
$4.66 + 0.20 - 1.22 = \qty{3.64}{g}$، أي \qty{3.22}{\percent} مقابل 3.19. \ce{K2O}:
$0.88 + 0.92 = \qty{1.80}{g}$، أي \qty{1.60}{\percent} مقابل 1.75. (ج) السيليكا والمغنيسيا، وهما اثنان من
الضوابط الثلاثة، تلائمان في حدود 0.03؛ أمّا البوتاس، الذي لم يُستعمل ضابطًا، فأدنى بمقدار 0.15، أي نحو
عُشر قيمته: فالمادة المُستوعَبة كان فيها بوتاس أكثر من متوسط
الصخر الدخيل، أو أُخذ مصهور منها غني بالبوتاس بالتفضيل.
\end{solution}

\begin{solution}{exo:b2:magmatic-differentiation:10}
(أ) $z = (0.2 + 1.5 - 1)/(0.2 - 1) = 0.7/(-0.8) = -0.875$؛ $F^{-z} = 0.7^{0.875} =
0.7319$؛ $C^*/C_0 = 0.2\times0.5/(0.2 - 1 + 1.5) = 0.1/0.7 = 0.1429$، ومنه $C/C_0 =
0.7319 + 0.1429\times(1 - 0.7319) = 0.7319 + 0.0383 = 0.770$. (ب) رايلي:
$0.7^{0.5} = 0.837$. \omterm{def:b2:magma-emplacement:magma}{الصهارة} في نموذج \omterm{def:b2:magmatic-differentiation:open}{الاستيعاب مع التبلور التجزيئي} فيها سترونشيوم أقل لسببين: فلكي تفقد
\qty{30}{\percent} من كتلتها وهي تكسب $r$ ضعف ما تبلوره يجب أن
تبلور $(1 - 0.7)/(1 - 0.2) = 0.375$ بدلًا من 0.30، والسترونشيوم الذي تكسبه
يأتي من صخر نصفها غنًى. (ج) $(0.7319\times0.7035 + 0.0383\times0.7200)/0.7702
= (0.5149 + 0.0276)/0.7702 = 0.70432$: ارتفعت النسبة بمقدار 0.0008 بينما
كان التبلور وحده سيبقيها عند 0.7035.
\end{solution}

\begin{solution}{exo:b2:magmatic-differentiation:11}
(أ) $1.1\times(1100 - 200) + 250 = 990 + 250 = \qty{1240}{J}$ لكل غرام. (ب)
$25.4\times1240 = \qty{31496}{J}$، قريبًا من \qty{31350}{J} عند ويلكوكس. (ج)
$31496/370 = \qty{85.1}{g}$ من البلاجيوكليز. (د) لم تُزل الموازنة الكتلية إلا $2.9 +
9.6 = \qty{12.5}{g}$ من البلّورات، وحرارتها الكامنة (نحو \qty{4600}{J}) تغطي أقل
من سُدس الحاجة؛ $r = 25.4/12.5 = 2.03$، أي أن الصخر المُستوعَب ضعف
البلّورات المتكوّنة، بينما تسمح الموازنة الحرارية بقيمة $r$ أدنى من 1 بكثير. \omterm{def:b2:magma-emplacement:magma}{فالصهارة} الثائرة
لا يمكن أن تكون قد وفّرت الحرارة بنفسها: فلا بدّ أن كتلة أكبر من \omterm{def:b2:magma-emplacement:magma}{الصهارة} قد تبلورت
وبردت في العمق، بعيدًا عن الأنظار.
\end{solution}

\begin{solution}{exo:b2:magmatic-differentiation:12}
(أ) المرحلة 1، $F$ من 1 إلى 0.8 مع $D = 5$: العامل $0.8^{4} = 0.4096$. المرحلة 2،
من 0.8 إلى 0.5 مع $D = 1.5$: يُطبَّق القانون على السائل المتبقّي، فيكون
العامل $(0.5/0.8)^{0.5} = 0.7906$. النيكل: $250\times0.4096\times0.7906 =
\qty{81.0}{ppm}$. (ب) $250\times0.5^{4} = \qty{15.6}{ppm}$، أي أدنى بخمس مرات مما ينبغي. (ج)
يفترض قانون رايلي أن $D$ ثابت؛ فحين تتغيّر المجموعة المتبلورة
يقفز $D$، ويجب أن يُستأنف القانون من تركيب السائل الذي بُلغ عند
التغيّر، مع قياس $F$ نسبةً إلى ذلك السائل.
\end{solution}

\begin{solution}{pb:b2:magmatic-differentiation:1}
\noindent\textbf{الجزء الأول.}

\textbf{1.} $\bar D_{\ce{Ni}} = 0.20\times5.0 + 0.50\times0 + 0.30\times2.0 = 1.60$.

\textbf{2.} $\bar D_{\ce{Zr}} = 0.20\times0.0007 + 0.50\times0.001 + 0.30\times0.195 =
0.00014 + 0.0005 + 0.0585 = 0.059$.

\textbf{3.} المنطقة العليا: $\bar D_{\ce{Ni}} = 0.45\times2.0 = 0.90$؛ $\bar D_{\ce{Zr}} =
0.55\times0.001 + 0.45\times0.195 = 0.088$.

\textbf{4.} النيكل متوافق في المنطقة الدنيا (1.60) ويصبح غير متوافق
قليلًا (0.90) حين يختفي الأوليفين؛ والزركونيوم غير متوافق في كلتيهما.

\textbf{5.} الكلينوبيروكسين: فهو يوفّر $0.0585$ من القيمة $0.059$ في المنطقة الدنيا، أي
كلها تقريبًا، لأنه المعدن الوحيد الذي معامله بعيد عن الصفر.

\textbf{6.} تتغيّر $(C_{\ce{Ni}}/C_{\ce{Zr}})$ بالعامل $F^{D_{\ce{Ni}} - D_{\ce{Zr}}}$؛
ولكل انتصاف، $0.5^{1.60 - 0.059} = 0.5^{1.541} = 0.344$.

\noindent\textbf{الجزء الثاني.}

\textbf{7.} $1.60\times150 = \qty{240}{ppm}$.

\textbf{8.} كان في السائل $190/1.60 = \qty{118.75}{ppm}$؛ $F = (118.75/150)^{1/0.6} =
0.7917^{1.667} = 0.68$.

\textbf{9.} $80\times0.68^{-0.941} = 80\times1.438 = \qty{115}{ppm}$.

\textbf{10.} التجزيئي: $150\times0.68^{0.6} = 150\times0.793 = \qty{119.0}{ppm}$؛
المتوازن: $150/(1.6 + 0.68\times(1 - 1.6)) = 150\times0.839 = \qty{125.8}{ppm}$.
والفرق \qty{6}{\percent} فقط بعد أن تبلور ثلث \omterm{def:b2:magma-emplacement:magma}{الصهارة}: فمن أجل
عنصر متوسط التوافق يبقى القانونان متقاربين في البداية.

\textbf{11.} $150\times0.55^{0.6} = \qty{104.8}{ppm}$.

\textbf{12.} تحته مباشرة: $1.60\times104.8 = \qty{167.7}{ppm}$؛ فوقه مباشرة:
$0.90\times104.8 = \qty{94.3}{ppm}$. فمقطع النيكل في المتراكمات ينخفض
فجأة بنسبة \qty{44}{\percent} عند المستوى الذي يختفي فيه الأوليفين من
الطبقات: درجة، لا منحنى أملس.

\textbf{13.} السائل: $104.8\times(0.40/0.55)^{0.90 - 1} = 104.8\times1.032 =
\qty{108.2}{ppm}$؛ \omterm{def:b2:magmatic-differentiation:cumulate}{والمتراكم} $0.90\times108.2 = \qty{97.4}{ppm}$. ومع $\bar D < 1$
يرتفع النيكل الآن ببطء في السائل وفي المتراكمات: فالمقطع، بعد
درجته، يصعد بلطف.

\noindent\textbf{الجزء الثالث.}

\textbf{14.} $x_{\mathrm{ol}}(m_{\mathrm{ol},j} - d_j) + x_{\mathrm{pl}}
(m_{\mathrm{pl},j} - d_j) + x_{\mathrm{cpx}}(m_{\mathrm{cpx},j} - d_j) = 100\,(p_j -
d_j)$.

\textbf{15.} \ce{MgO}: $37.16\,x_{\mathrm{ol}} - 4.9\,x_{\mathrm{pl}} +
11.1\,x_{\mathrm{cpx}} = 410$؛ \ce{CaO}: $-7.3\,x_{\mathrm{ol}} +
7.06\,x_{\mathrm{pl}} + 13.2\,x_{\mathrm{cpx}} = 420$؛ \ce{Al2O3}:
$-14.1\,x_{\mathrm{ol}} + 17.6\,x_{\mathrm{pl}} - 11.1\,x_{\mathrm{cpx}} = 190$.

\textbf{16.} بالحلّ (بالحذف أو بقاعدة كرامر): $x_{\mathrm{ol}} = 9.03$،
$x_{\mathrm{pl}} = 30.84$، $x_{\mathrm{cpx}} = 20.32$ غرامًا لكل \qty{100}{g}.

\textbf{17.} البلّورات $9.03 + 30.84 + 20.32 = \qty{60.19}{g}$؛ والسائل
\qty{39.81}{g}: $F = 0.398$.

\textbf{18.} $(4950 - 9.03\times39.19 - 30.84\times50.54 - 20.32\times51.5)/39.81 =
(4950 - 353.9 - 1558.7 - 1046.5)/39.81 = 1990.9/39.81 = \qty{50.0}{\percent}$، كما في
\omterm{def:b2:magma-emplacement:bodies}{القاطع}: فالأكسيد غير المستعمل ضابطًا يُستعاد.

\textbf{19.} البلّورات المُزالة هي $9.03/60.19 = \qty{15}{\percent}$ من الأوليفين،
و\qty{51}{\percent} من البلاجيوكليز و\qty{34}{\percent} من الكلينوبيروكسين، بين
التركيبين المعدنيين للمنطقتين. وهي تطابق المجموعة الدنيا متبلورةً من $F = 1$ إلى
0.55 ($0.45\times20 = 9.0$، $0.45\times50 = 22.5$، $0.45\times30 = 13.5$) تليها
المجموعة العليا من 0.55 إلى 0.40 ($0.15\times55 = 8.25$، $0.15\times45 = 6.75$): أي 9.0
و30.75 و20.25 غرامًا.

\textbf{20.} $0.30/0.75 = 0.40$.

\noindent\textbf{الجزء الرابع.}

\textbf{21.} $0.55^{0.059 - 1} = 0.55^{-0.941} = 1.755$؛ $80\times1.755 =
\qty{140.4}{ppm}$.

\textbf{22.} $188/140.4 = 1.339 = (F/0.55)^{0.088 - 1}$، ومنه $F = 0.55\times1.339^{-1/
0.912} = 0.55\times0.726 = 0.399$.

\textbf{23.} $(188/80)^{1/(0.07 - 1)} = 2.35^{-1.075} = 0.399$. فللعناصر غير المتوافقة
$C_L/C_0 \approx 1/F$ أيًّا كانت قيمة $D$ الصغيرة: فالجواب لا يكاد يتوقّف على
تفاصيل $D$، وهذا ما يجعل العناصر غير المتوافقة أفضل مقاييس $F$.

\textbf{24.} البلاجيوكليز، الذي لا يحوي حديدًا، يشكّل نصف البلّورات،
والبيروكسين والأوليفين معًا يحويان من الحديد أقل مما يحويان من المغنيسيا: فيزداد السائل
غنًى بالحديد وفقرًا بالمغنيسيا دون أن يكسب سيليكا. وهذا الإثراء
بالحديد عند سيليكا شبه ثابتة هو علامة المتتالية الثولييتية في
الفصل التالي.

\textbf{25.} البلّورات المتراكمة النقية تحوي $0.088\times188 = \qty{16.5}{ppm}$. ومع
كسر كتلي $t$ من السائل المحبوس، $16.5(1 - t) + 188\,t = 47$، ومنه $t = (47 -
16.5)/(188 - 16.5) = 30.5/171.5 = 0.18$: أي نحو \qty{18}{\percent} من السائل المحبوس.

\textbf{26.} التقديرات الثلاثة، العناصر الرئيسية (0.398) والبوتاس (0.40)
والزركونيوم (0.399)، متّفقة: فحين تبلورت الطبقة العليا كان
\textbf{\qty[text-series-to-math=true]{40}{\percent}} من \omterm{def:b2:magma-emplacement:magma}{الصهارة} الأصلية لا يزال
سائلًا.
\end{solution}
